\section{11. DFS}

\subsection{Домашнее задание}
\begin{enumerate}
  \item
    Пусть дано дерево $T = \langle V, E \rangle$. Для каждого ребра вычислить, сколько
    простых путей проходит через него. $\O(|V| + |E|)$.

    Решение:

    \begin{itemize}
      \item Каждое ребро разбивает граф на 2 компоненты связности. В 1-ой компоненте мы можем выбрать, где начинаем путь, а во 2-ой -- где заканчиваем.

      \item Тогда в случае подвешенного дерева на $N$ вершинах количество простых путей, проходящих через ребро, будет равно $S * (N - S)$, где
      \newline $S$ -- это количество способов выбрать точку начала (то же самое, что размер 1-ой компоненты связности), 
      \newline $(N - S)$ -- это количество способов выбрать точку конца (то же самое, что размер 2-ой компоненты связности)

      \item Получается, надо для каждого ребра посчитать размер одной из его компонент связности. Для этого можно запустить обычный DFS. 
      Однако это неоптимально -- сложность будет $\O(|E| \cdot |V|) = \O(|V|^2)$, так как в дереве $|E| = |V| - 1$.

      \item Оптимальнее будет запустить один DFS и посчитать размер поддеревьев всех вершин за $\O(|V|)$:
      \newline при выходе из вершины $v$ её размер равен $1 + \sum_{u \in children(v)} size(u)$

      \item Тогда для каждого ребра количество простых путей вычисляется за $\O(1)$
      \item Сложность по времени: $\O(|V| + |E|) = \O(|V|)$.
      \item Сложность по памяти: $\O(|V|)$ — массив для хранения размеров поддеревьев

    \end{itemize}

  \item
    Корневое дерево $T$ на $V$ вершинах задается массивом из $V$ элементов. 
    Все вершины пронумерованы. Для каждой вершины в массиве указан его родитель.
    Для корня $r$ значение в массиве равно -1. 
    Требуется определить, как будет выглядеть новое представление 
    дерева, если корень $r$ сменить на корень $q$. Разрешается использовать
    $\O(1)$ дополнительной памяти. Менять массив можно. Время $\O(V)$.

    Решение:

    \begin{itemize}
      \item При смене корня с $r$ на $q$ нужно развернуть рёбра только на пути от $q \leadsto r$.
      Все остальные вершины дерева уже корректно указывают на своих родителей.

      \item Для этого достаточно держать 3 указателя:
      \begin{itemize}
        \item на текущую вершину
        \item на предыдущую вершину на пути $q \leadsto r$ (она станем новым родителем текущей вершины)
        \item на старого родителя текущей вершины (для перемещения по дереву)
      \end{itemize}

      \item На каждом шаге сохраняем старого родителя, меняем родителя на предыдущую вершину пути, переходим дальше.

      \item Сложность по времени: $\O(|V|)$ — длина пути от $q \leadsto r$ не превышает $|V|$

      \item Сложность по памяти: $\O(1)$ — только 3 переменные.
    \end{itemize}


  \item
    За $\O(|V| + |E|)$ найдите в \textbf{неорграфе}:
    \begin{enumerate}
      \item все \textbf{ребра}, через которые проходит любой путь $a \leadsto b$.
      
      Решение:

      \begin{itemize}
        \item Допустим, мы нашли через DFS какой-то путь $a \leadsto b$. Все интересующие нас ребра точно будут на этом пути.
        
        Если на этом пути есть ребро, которое при удалении разобьет граф на 2 компоненты связности,
        то это значит, что мы нашли мост. И оно же является ребром, который лежит на всех путях $a \leadsto b$

        \item Получается, что для решения задачи нужно найти все мосты на пути $a \leadsto b$.

        \item Поиск мостов за $\O(|V| + |E|)$:

        Пусть $enter(x)$ — время входа в вершину $v$ при DFS.
        
        Пусть $ret(x)$ — минимальное время входа среди вершин, достижимых из поддерева с корнем $x$:
        $$ret(x) = \min\begin{cases}
          enter(x) \\
          \min_{(u, x) \in up} \{ enter(u) \}, & \text{где } up \text{ — множество обратных ребер} \\
          \min_{z \in V} \{ ret(z) \}, & \text{где } z \text{ — потомок } x
        \end{cases}$$

        Ребро $(x, y)$ является мостом $\Leftrightarrow ret(y) > enter(x)$. Если это условие выполняется,
        то из поддерева $y$ нельзя попасть в $x$ или выше без ребра $(x, y)$

        \item Финальный алгоритм:
        \begin{enumerate}
          \item Найти путь $a \leadsto b$ через DFS
          \item Найти все мосты графа за $\O(|V| + |E|)$
          \item Вернуть все мосты, принадлежащие пути $a \leadsto b$
        \end{enumerate}

        \item Сложность по времени: $\O(|V| + |E|)$
      \end{itemize}

      \item все \textbf{вершины}, через которые проходит любой путь $a \leadsto b$.
      
      Решение:

      \begin{itemize}

        \item Рассуждение аналогичное: берем путь $a \leadsto b$. Все интересующие нас вершины будут на этом пути.
        
        Если на этом пути есть вершина $x$, которая при удалении разобъет граф на $\geq 2$ компоненты связности,
        то значит $x$ -- это точка сочленения. И она же является вершиной, которая лежит на всех путях $a \leadsto b$.

        \item Поиск точек сочленения за $\O(|V| + |E|)$:

        Вершина $x$ является точкой сочленения $\Leftrightarrow$
        \begin{itemize}
          \item $x$ — корень дерева и имеет $\geq 2$ детей, или
          \item $x$ — не корень и существует ребёнок $y$ такой, что $ret(y) \geq enter(x)$
        \end{itemize}

        \item Финальный алгоритм:
        \begin{enumerate}
          \item Найти путь $a \leadsto b$ через DFS
          \item Найти все точки сочленения графа за $\O(|V| + |E|)$
          \item Вернуть точки сочленения, лежащие на пути $a \leadsto b$ (кроме самих $a$ и $b$)
        \end{enumerate}

        \item Сложность по времени: $\O(|V| + |E|)$
      
      \end{itemize}

    \end{enumerate}

  \item
    Дано дерево $T = \langle V, E \rangle$.
    \begin{enumerate}
      \item
        Найдите центроид дерева --- вершину $c$, для которой минимальна \\
        $\sum_{v \in V} dist(v, c)$. $\O(|V|)$.

        Решение:

        \begin{itemize}
          
          \item Возьмем произвольную вершину $x$ и подвесим за нее дерево

          \item Обойдем все дерево через DFS и посчитаем формулу центроида из условия. Получили какое-то значение $d$

          \item Вполне возможно, что это не является ответом на задачу. 
          Однако пересчитывать формулу для каждой вершины неоптимально - будет $\O(|V|^2)$.
          
          \item Вместо этого мы можем сделать следующее:
          \begin{itemize}
            \item Пусть $y$ -- один из потомков $x$

            \item Если мы сдвинемся в вершину $y$, то все расстояния в поддереве $y$ уменьшаться на 1,
            а в остальной части дерева -- увеличатся на 1

            \item Тогда получается формула пересчета: $d - S + (N - S)$, где
            \newline $S$ -- размер поддерева с вершиной $y$,
            \newline $(N - S)$ - размер остальной части дерева.
          \end{itemize}

          \item Таким образом обходим все поддерево, пересчитываем значение по новой формуле и получаем ответ.
          
          \item Сложность по времени: $\O(|V|)$
        
        \end{itemize}

    \end{enumerate}

\end{enumerate}
